\section{Introduction}
\label{sec:intro}

\texttt{dc:title} values in the Deutsche Digitale Bibliothek are inconsistent:
some contain only a proper title, while many follow ISBD conventions that pack
subtitle, statement of responsibility~--- including role markers such as
\textit{hrsg.\ v.}~--- and imprint details into the same field.
The properties that carry person information offer no improvement:
\texttt{dc:contributor} conflates editors, translators, and publishers without
role distinction, and \texttt{dc:creator} is absent in the majority of records.
None of these field values can be used for Work-level authority reconciliation,
nor do they indicate whether a named person authored, edited, or translated a work.
The paper addresses this with a silver labeling pipeline that leverages the ISBD
delimiter structure present in 7.56\% of 9.2M German-language records to extract
\texttt{TITLE}, \texttt{OTHER\_TITLE}, and \texttt{PERSON} spans, and evaluates
five systems~--- from a zero-shot baseline (NuNER~Zero, $F_1 = 0.000$) to
fine-tuned GLiNER and XLM-R~large~--- against a 395-record gold set stratified
by era and silver tier.

% TODO: contribution list, section roadmap.

\endinput
%% Everything below is storyline planning notes — not compiled.
% ============================================================
% STORYLINE — planning notes, not for compilation
% Based on: paper-framing-jcdl.md §2, spiel_ner.md §0--§3,
%           jcdl-plan-202606.md §1, assessment rounds §5--§6
% ============================================================

% -----------------------------------------------------------------------
\subsection*{0. Narrative arc}
% Three moves, each motivating the next.
\begin{enumerate}
    \item Bibliographic catalog title strings are semi-structured: ISBD punctuation
          encodes structure in 7.56\% of records; the remaining 92.44\% are
          free-text. Both halves fail for general NER, but for different reasons
          that the three failure dimensions capture.
    \item That same ISBD structure is a reliable labeling signal where it exists.
          The pipeline contribution exploits the thing that breaks general NER.
    \item The 92.44\% without structural signal requires a different strategy:
          that is future work. This paper establishes the silver labeling baseline
          and evaluation protocol for the tractable subset.
\end{enumerate}
% Opening paragraph:
\begin{quote}
    Bibliographic catalog strings are semi-structured.
    In 7.56\% of 9.21 million German-language records in the Deutsche Digitale
    Bibliothek, ISBD punctuation encodes title, subtitle, and authorship in a
    fixed grammar --- ` : ' separating title from subtitle, ` / ' separating
    title block from statement of responsibility.
    In the remaining 92.44\%, the same \texttt{dc:title} field carries a free-text
    string with no recoverable boundary markers.
    Both cases defeat off-the-shelf named entity recognition, but along three
    independent dimensions: the target label set does not exist in standard
    training corpora, the domain is GLAM bibliographic rather than newswire,
    and the register alternates between formulaic punctuation grammar and
    unstructured prose.
    We address the ISBD-covered stratum with a silver labeling pipeline that
    treats the existing punctuation grammar as a labeling signal, a 395-record
    stratified gold evaluation set, and a five-system comparison establishing
    zero-shot and fine-tuned baselines for Work-level entity extraction from
    historical German library catalogs.
\end{quote}
% NAMING RULE: "DDB" / "Deutsche Digitale Bibliothek" named once above.
% All subsequent references: "9.21M records", "the corpus", "the catalog strings".
%
% DOUBLE-BLIND: cc2llm \cite{tan-etal-2024-card} → ``[Author, 2024]''.
% Self-citation rule: GeMeA ISWC, babel-ddb, mocho → ``[Author, 2026]''.

% -----------------------------------------------------------------------
\subsection*{1. The domain: why this corpus, why now}
\begin{itemize}
    \item 9.21M German-language DDB catalog records; libraries sector;
          \texttt{dc:title} present in 100\% of records --- often the only
          structured metadata field carrying entity-level information
    \item Title field is \textbf{semi-structured}: ISBD punctuation present in
          7.56\% (tier-1: 689K records, tier-2: 7K records); remaining 92.44\%
          (tier-0) are free-text --- both require NER but respond to different strategies
    \item Downstream motivation: GeMeA Phase 0 --- Work-level entity extraction
          is the prerequisite for GND authority linking; cannot be skipped
    \item Prior work \cite{tan-etal-2024-card} established NER baselines on a
          108K subset; present work scales to 9.21M and adds silver labeling +
          fine-tuning on the ISBD-covered stratum
          % cc2llm: cite as [Author, 2024] in submission
\end{itemize}

% -----------------------------------------------------------------------
\subsection*{2. Three failure dimensions of off-the-shelf NER}
% Structural backbone of §4 (Silver and Gold Dataset).
% Three independent dimensions — each alone would motivate domain adaptation;
% together they make zero-shot transfer from standard corpora unviable.
\begin{itemize}
    \item \textbf{Wrong label set}: CoNLL/OntoNotes PER, ORG, LOC have no overlap
          with TITLE, OTHER\_TITLE, PERSON as defined for FRBR Work identification;
          zero positive training signal for the target types
    \item \textbf{Wrong domain}: NER training corpora are newswire, Wikipedia, web text;
          GLAM bibliographic records are out-of-distribution across vocabulary,
          named entity density, and entity boundary conventions
    \item \textbf{Wrong register}: the ISBD-covered stratum uses punctuation as
          structural markers rather than syntactic cues; the free-text stratum
          uses long title-page transcriptions with pre-title author placement and
          Early Modern orthography (pre-1750) --- both defeat prose-trained models,
          but differently
    \item Evidence: NuNER Zero F1 = 0.000 on tier-2 records across all label types
          and all prompt variants (SR-09, 2026-03-27)
\end{itemize}

% -----------------------------------------------------------------------
\subsection*{3. ISBD as the labeling signal}
% The thing that causes failure in general NER enables silver labeling.
% This is the core originality of the approach.
\begin{itemize}
    \item ISBD punctuation is a deterministic segmentation signal where it exists:
          ` : ' $\Rightarrow$ TITLE / OTHER\_TITLE boundary;
          ` / ' $\Rightarrow$ statement of responsibility $\Rightarrow$ PERSON candidate
    \item Tiered confidence: tier-2 (full structural pattern) $>$ tier-1
          (heuristic, partial) $>$ tier-0 (no signal; out of scope)
    \item Tier-0 scope statement (required per Assessment \#2 §6.4):
          ``We focus on the 7.56\% of records with ISBD structural signal.
          Tier-0 records (92.44\%) require a different strategy and are deferred.''
    \item Fine-tuning on silver tests whether operationalised domain knowledge ---
          ISBD rules encoded as labeling heuristics --- constitutes a sufficient
          training signal; that is the validation question, not silver labeling
          as a general method
          % Do NOT frame as "a general silver labeling method" (Assessment #2 §6.5)
\end{itemize}

% -----------------------------------------------------------------------
\subsection*{4. Contributions}
\begin{enumerate}
    \item \textbf{ISBD-guided silver labeling pipeline} with tiered confidence
          (tier-0/1/2); 689K silver spans from tier-1, 7K from tier-2;
          char-offset triples in GLiNER/BIO-compatible format
    \item \textbf{395-record stratified gold evaluation set} with human-human IAA
          (frankp $\cap$ maria: 131 records; frankp $\cap$ gundi: 82 records);
          era $\times$ silver tier stratification; Doccano JSONL format
    \item \textbf{Five-system comparison}: NuNER Zero, GLiNER-multi zero-shot,
          GLiNER-multi few-shot (k-ablation), GLiNER-multi fine-tuned on silver,
          XLM-R large fine-tuned on silver (BIO); k-fold cross-validation;
          per-stratum F1 $\pm$ std
    \item \textbf{Dataset and eval-script release}: silver spans + gold set +
          evaluation code (CC BY 4.0; anonymous.4open.science/XXX)
\end{enumerate}
% WORDING NOTE (Assessment #2 §6.3):
% Never use "ceiling" — use "zero-shot and fine-tuned baselines" or
% "baseline performance characterisation".
%
% BibNER-WERK PHASE 1 FOOTNOTE (Assessment #2 §6.6): add near first mention of "BibNER-Werk":
% "Phase 1 covers Work-level entities (TITLE, OTHER\_TITLE, PERSON).
%  Expression and Manifestation levels are future phases."

% -----------------------------------------------------------------------
\subsection*{5. Section roadmap}
\begin{quote}
    Section~\ref{sec:related} situates the work in historical NER and
    bibliographic NLP; Section~\ref{sec:corpus} describes the corpus and task;
    Section~\ref{sec:silverandgold} presents the silver labeling pipeline and
    gold evaluation set; Section~\ref{sec:setup} details the experimental setup;
    Section~\ref{sec:results} reports results; Section~\ref{sec:conclusion} concludes.
\end{quote}

% -----------------------------------------------------------------------
\subsection*{Notes and open questions}
\begin{itemize}
    \item Opening sentence ``Bibliographic catalog strings are semi-structured''
          --- check whether this is the right entry point or whether
          the scale claim (9.21M) should come first.
    \item First use of ``BibNER-Werk'' needs Phase 1 footnote --- add after
          contribution list or in abstract.
    \item \texttt{fig\_title\_lengths.png} (sr10) --- belongs in
          §\ref{sec:corpus}, not §1.
    \item Confirm GeMeA citation treatment for double-blind ---
          see \texttt{paper-framing-jcdl.md §0.1}.
    \item Re-read \texttt{spiel\_ner.md §0--§3} before drafting ---
          maps nearly verbatim to §§1--4.
\end{itemize}
%% end storyline
